\documentclass[11pt, a4paper]{article}

% --- PACCHETTI ---
\usepackage[utf8]{inputenc}
\usepackage[italian]{babel}
\usepackage[margin=2.5cm]{geometry} % Imposta i margini della pagina
\usepackage{graphicx}               % Per inserire il logo
\usepackage{tabularx}               % Per tabelle larghe quanto la pagina
\usepackage[table]{xcolor}          % Per colorare le celle delle tabelle
\usepackage[hidelinks]{hyperref}    % Per stile link cliccabili
\usepackage{lmodern}      % Rende i font scalabili all'infinito
\usepackage{anyfontsize}  % Permette di usare dimensioni personalizzate
\usepackage{xurl}         % Permette ai link lunghi di non uscire dai margini
\urlstyle{same}           % Usa lo stesso font del documento

\begin{document}

% ==============================================================================
% FRONTESPIZIO (PAGINA 1)
% ==============================================================================
\begin{titlepage}
    \centering
    \vspace*{0.5cm}

    % LOGHI
    \includegraphics[width=0.45\textwidth]{assets/logo_unipd.jpg}
    \hfill
    \includegraphics[width=0.45\textwidth]{assets/logo.png}\par
    \vspace{0.5cm}

    % NOME DEL GRUPPO
    {\fontsize{45}{50}\sffamily\selectfont \textbf{On-By-One}\par} 
    \vspace{1cm}

    % TITOLO DEL DOCUMENTO
    {\huge \bfseries Analisi dei requisiti \par}
    
    \vspace{1.5cm}

    % MEMBRI DEL GRUPPO
    {\large
    \textbf{On-By-One}:\par
    \vspace{0.3cm}
    \begin{tabular}{l l}
        Sara Ballardin & 2101048 \\
        Shuhui Chen & 2138001 \\
        Simone Giuliano & 2145591 \\
        Alvise Greggio & 2141004 \\
        Daniele Nuzzaci & 2139039 \\
        Federico Temporin & 2146764 \\
        Tommaso Zorzetto & 2148629 \\
    \end{tabular}
    \par}

    \vfill

    % TABELLA INFORMAZIONI DOCUMENTO E DESTINATARI
    \begin{table}[h]
        \centering
        \arrayrulecolor{gray!50}
        \renewcommand{\arraystretch}{1.5}
        \begin{tabularx}{\textwidth}{|c|c|c|X|}
            \hline
            \rowcolor{gray!10} \multicolumn{4}{|c|}{\textbf{Informazioni documento}} \\
            \hline
            \rowcolor{gray!10} \textbf{Versione} & \textbf{Data} & \textbf{Stato} & \multicolumn{1}{c|}{\textbf{Destinatari}} \\
            \hline
            - & - & Bozza & 
            On-By-One \newline 
            Prof. Tullio Vardanega \newline 
            Prof. Riccardo Cardin \newline 
            Vimar S.P.A. \\
            \hline
        \end{tabularx}
        \arrayrulecolor{black}
    \end{table}
    
    \vspace{0.5cm}
    
    % CONTATTI
    {\large Contatti: \href{mailto:onbyone.swe@gmail.com}{onbyone.swe@gmail.com} }
    
\end{titlepage}

\newpage


% ==============================================================================
% REGISTRO DELLE MODIFICHE (PAGINA 2)
% ==============================================================================
\section*{Registro delle modifiche}

\vspace{1cm}
\begin{center}
    \arrayrulecolor{gray!50}
    \renewcommand{\arraystretch}{1.5}
    \begin{tabularx}{\textwidth}{|c|c|c|X|c|}
        \hline
        \rowcolor{gray!10} 
        \multicolumn{1}{|c|}{\textbf{Versione}} & 
        \multicolumn{1}{c|}{\textbf{Data}} & 
        \multicolumn{1}{c|}{\textbf{Autore}} & 
        \multicolumn{1}{c|}{\textbf{Descrizione}} & 
        \multicolumn{1}{c|}{\textbf{Verificatore}} \\
        \hline
        - & - & - & - & - \\
        \hline
        0.1.1 & 2026/11/13 & Sara Ballardin & Stesura dell'introduzione & - \\
        \hline
        0.1.0 & 2026/11/09 & Sara Ballardin & Definizione struttura del documento & - \\
        \hline
    \end{tabularx}
    \arrayrulecolor{black}
\end{center}
\vfill
\newpage


% ==============================================================================
% INDICE (PAGINA 3)
% ==============================================================================
\tableofcontents
\newpage


% ==============================================================================
% 1. INTRODUZIONE
% ==============================================================================
\section{Introduzione}

\subsection{Scopo del documento}
Il documento di \textbf{Analisi dei requisiti} raccoglie e descrive le funzionalità e i vincoli del progetto \textbf{AURA} (AI for Understanding, Reasoning and Automation), proposto da Vimar S.P.A. e sviluppato dal gruppo On-By-One.

\noindent L'obiettivo è trasformare le richieste del capitolato e quanto emerso negli incontri con l'azienda in un elenco chiaro e verificabile di requisiti software. Partendo dallo studio degli attori e dei casi d'uso in UML, il documento definisce cosa deve fare il sistema per soddisfare le esigenze dell'utente, lasciando le scelte su come realizzarlo alla successiva fase di progettazione.

\noindent Il documento rappresenta il punto di riferimento per tutti i soggetti coinvolti nel progetto:
\begin{itemize}
    \item \textbf{Proponente (Vimar S.P.A.) e committenti (Prof.\ Tullio Vardanega e Prof.\ Riccardo Cardin):} per verificare che le richieste iniziali siano state comprese correttamente e nella loro interezza;
    \item \textbf{Progettisti e programmatori:} come guida per definire l'architettura del sistema e sviluppare le funzionalità previste;
    \item \textbf{Verificatori:} come base per scrivere i test e controllare che il prodotto realizzato rispetti tutti i requisiti fissati.
\end{itemize}

\subsection{Glossario}
Per evitare ambiguità e facilitare la lettura di termini tecnici, acronimi e parole con un significato specifico nel progetto, il gruppo ha redatto un documento interno dedicato chiamato \textbf{\href{https://on-by-one-team.github.io/On-By-One.github.io/docs/pdf/glossario/glossario.pdf}{Glossario}}, che raccoglie tutte le relative definizioni.

\subsection{Riferimenti}
Di seguito sono elencati i documenti vincolanti di progetto, le regole interne del gruppo e le fonti informative (standard internazionali e materiale del corso) consultate per redigere questo documento.

\subsubsection{Riferimenti normativi}
\begin{itemize}
    \item Capitolato d'appalto AURA (Vimar S.P.A.) \\
    \url{https://www.math.unipd.it/~tullio/IS-1/2026/Progetto/C1-AURA.pdf} \\
    Ultimo accesso: 13 novembre 2026.

    \item Regolamento del progetto didattico \\
    \url{https://www.math.unipd.it/~tullio/IS-1/2026/Dispense/PD1.pdf} \\
    Ultimo accesso: 13 novembre 2026.

    \item Norme di progetto (documento interno del gruppo On-By-One) \\
    \url{https://on-by-one-team.github.io/On-By-One.github.io/docs/pdf/norme-di-progetto-documentazione.pdf} \\
    Ultimo accesso: 13 novembre 2026.
\end{itemize}

\subsubsection{Riferimenti informativi}
\begin{itemize}
    \item Analisi dei requisiti (Prof.\ Tullio Vardanega) \\
    \url{https://www.math.unipd.it/~tullio/IS-1/2026/Dispense/T05.pdf} \\
    Ultimo accesso: 13 novembre 2026.

    \item Diagrammi dei casi d'uso (Prof.\ Riccardo Cardin) \\
    \url{https://www.math.unipd.it/~rcardin/swea/2022/Diagrammi%20Use%20Case.pdf} \\
    Ultimo accesso: 13 novembre 2026.
\end{itemize}

\newpage


% ==============================================================================
% 2. DESCRIZIONE DEL PRODOTTO
% ==============================================================================
\section{Descrizione del prodotto}

\subsection{Scopo e obiettivi del prodotto}
Descrivere lo scopo del prodotto e i traguardi generali che il sistema deve raggiungere, formalizzando le aspettative espresse dalla proponente.

\subsection{Contesto d'uso e prospettiva del prodotto}
Inquadrare il sistema nel suo ambiente operativo: interazioni con gli utenti, componenti che lo costituiscono e servizi/piattaforme esterne con cui scambia informazioni.

\subsection{Funzionalità principali}
Presentare una panoramica delle macro-funzionalità fornite dal sistema, senza anticipare il livello di dettaglio dei requisiti software.

\subsection{Caratteristiche degli utenti}
Descrivere i profili di utilizzatori umani e di sistemi esterni che si interfacciano con il prodotto, evidenziandone competenze tecniche attese, mansioni e obiettivi.

\newpage


% ==============================================================================
% 3. CASI D'USO (USE CASES)
% ==============================================================================
\section{Casi d'uso}

\subsection{Introduzione ai casi d'uso}
Spiegare cos'è un caso d'uso, la convenzione di codifica adottata (UC, sotto-casi d'uso) e come sono letti i diagrammi UML.
 
\subsection{Formato di specifica}
Definire i campi usati per descrivere ogni caso d'uso: attori, pre-condizioni, post-condizioni, scenario principale, scenari alternativi, inclusioni, estensioni e generalizzazioni.
 
\subsection{Classificazione degli attori}
\subsubsection{Attori principali}
Elencare e descrivere gli attori che avviano proattivamente le interazioni con il sistema per raggiungere un obiettivo di business, allegando il diagramma generale delle relazioni tra attori.
 
\subsubsection{Attori secondari}
Elencare e descrivere i sistemi, le API o i servizi esterni che prendono parte all'esecuzione dei casi d'uso senza avviarli direttamente.

\subsection{Elenco dei casi d'uso}
Raccogliere i casi d'uso raggruppati per attore principale o per area funzionale. Per ognuno inserire diagramma UML e descrizione secondo il formato definito sopra.

\newpage


% ==============================================================================
% 4. REQUISITI
% ==============================================================================
\section{Requisiti}

\subsection{Classificazione e codifica dei requisiti}
Descrivere la convenzione di codifica dei requisiti (tipologia, numero progressivo, priorità obbligatorio/desiderabile/opzionale) e le fonti da cui derivano (casi d'uso, capitolato, verbali degli incontri con la proponente).

\subsection{Requisiti funzionali}
Requisiti che descrivono cosa il sistema deve fare (codice identificativo, descrizione e fonte (caso d'uso)).

\subsection{Requisiti di qualità}
Requisiti sulla qualità di prodotto e di processo (testing, copertura, documentazione da consegnare, manutenibilità).

\subsection{Requisiti di vincolo}
Requisiti che impongono limiti alle scelte progettuali e tecnologiche.

\subsection{Requisiti prestazionali}
Requisiti misurabili su tempi di risposta, capacità e scalabilità, se presenti.

\newpage


% ==============================================================================
% 5. TRACCIAMENTO DEI REQUISITI
% ==============================================================================
\section{Tracciamento}

\subsection{Tracciamento fonti -- requisiti}
Tabella che mappa ciascuna fonte formale (capitolato d'appalto, verbali di riunione esterni, casi d'uso) ai requisiti che ne sono scaturiti.

\subsection{Tracciamento casi d'uso -- requisiti}
Tabella di correlazione che associa ogni singolo caso d'uso ai relativi requisiti funzionali soddisfatti.


% ==============================================================================
% 6. RIEPILOGO
% ==============================================================================
\section{Riepilogo dei requisiti}
Tabella riassuntiva con il conteggio complessivo dei requisiti per tipologia (funzionali, qualitativi, di vincolo, prestazionali) e per grado di priorità (obbligatori, desiderabili, opzionali).

\end{document}